% ============================================================
% Chapter 4 --- Biased Beliefs
% Converted from Slides/Lecture4_LaTeX/Lecture4_Beliefs.tex
% Figures reproduced in figures/ (original slide assets).
% Synced with revised deck 2026-09-25: individualism, DOX, SPF
% dispersion, resale-option, forecast-dispersion, beta/alpha and
% Hong--Sraer figures; two "My Research" subsections added.
% ============================================================
\chapter{Biased Beliefs}
\label{ch:beliefs}

The benchmark model of investor beliefs is Bayes' rule. Agents hold prior probabilities over states of the world, observe signals whose likelihoods they know, and update to posterior probabilities by Bayes' formula. Combined with the expected-utility benchmark of Chapter~\ref{ch:preferences}, rational Bayes--expected-utility agents trading in frictionless markets deliver the predictions reviewed in Chapter~\ref{ch:intro}---prices track rational forecasts of cash flows, expected returns in the cross-section line up with risk, and investors hold diversified portfolios and trade only on information. Like the expected-utility benchmark, the Bayes benchmark is the wrong descriptor of how investors actually form beliefs, and the descriptive gaps are large enough to matter for prices. Laboratory experiments, survey data, and field evidence document systematic departures from Bayes' rule---representativeness, base-rate neglect, sample-size neglect, overconfidence, over-extrapolation, conservatism, anchoring, belief perseverance---and these departures show up in trading behavior, in expectations surveys, and in the cross-sectional and time-series properties of asset returns.

This chapter develops that evidence. The first half (Sections~\ref{sec:representativeness}--\ref{sec:otherbiases}) is the \emph{psychology of beliefs}: a survey of the cognitive biases that distort investors' probability assessments, together with the asset pricing applications each bias has been claimed to explain. The second half (Sections~\ref{sec:hetsources}--\ref{sec:causal}) is the \emph{heterogeneous beliefs} literature: rather than assuming a single representative investor with biased beliefs, the second half dispenses with the representative assumption and asks what cross-sectional variation in beliefs does to prices---especially when combined with short-sale constraints. Section~\ref{sec:stakes} asks whether people ever learn, Section~\ref{sec:beliefscontested} catalogues the limits of the biased-beliefs program, and Section~\ref{sec:beliefsahead} closes the chapter.

% ============================================================
\section{Representativeness}
\label{sec:representativeness}

\begin{definition}[Representativeness heuristic]
\label{def:representativeness}
A person judging the probability that object $A$ belongs to class $B$, or that event $A$ was generated by process $B$, evaluates the probability by the degree to which $A$ resembles $B$---the \emph{degree to which $A$ reflects the essential characteristics of $B$} \citep{kahneman1974}.
\end{definition}

The representativeness heuristic is useful in many settings: the prediction that a small, dense, red fruit is more likely an apple than a cherry is exactly the kind of inference the heuristic supports. The bias enters when the resemblance judgment is made \emph{without} taking into account base rates and sample sizes, the two pieces of statistical information a Bayesian calculation would weight heavily. \citet{kahneman1974} document two robust violations.

\subsection{Base-Rate Neglect}

The standard demonstration is the ``Steve'' problem. Subjects read a personality sketch of a fictional person---very shy and withdrawn, invariably helpful, but with little interest in people or in the world of reality; a meek and tidy soul with a need for order and structure and a passion for detail---and are asked whether Steve is more likely to be a librarian or a lawyer. The modal answer is librarian, because the description resembles the librarian stereotype. But there are many more lawyers than librarians in the U.S., and a Bayesian would let that base rate dominate. The question asks about the posterior odds on the left of
\[
\frac{p(\text{librarian}\mid\text{data})}{p(\text{lawyer}\mid\text{data})}
=
\frac{p(\text{data}\mid\text{librarian})}{p(\text{data}\mid\text{lawyer})}
\cdot
\frac{p(\text{librarian})}{p(\text{lawyer})},
\]
and representativeness judges answer using only the likelihood ratio, the first term on the right. The prior odds are dropped, and with them the information that should dominate the answer. Base-rate neglect is the neglect of the second factor in a Bayes calculation the judge never performs.

\subsection{Sample-Size Neglect}

The second violation is the failure to weight evidence by the sample size from which it is drawn. A sequence of three heads from a fair coin is taken as strong evidence that the coin is biased toward heads, even though the posterior probability of bias after three flips is still close to the prior. The same judgment applied to a sequence of thirty heads would be more defensible. \citet{gilovich1985} document the consequence in basketball: fans and players judge a player who has made a few shots in a row to be ``hot,'' and adjust their expectations of his next shot upward. The empirical evidence in large samples is that the hot hand is largely a myth---shot-level success rates are essentially uncorrelated with the recent history of makes and misses---although \citet{millersanjurjo2018} later identified a selection bias in the standard tests that works against finding streaks, softening the original verdict. \citet{gilovich1985} read the phenomenon as representativeness: fans expect a short streak of makes to be representative of a player's true ability, and they expect the next shot to reflect that true ability. In finance, the same bias is the source of over-extrapolation, to which Section~\ref{sec:overextrapolation} returns.

% ============================================================
\section{Overconfidence}
\label{sec:overconfidence}

\begin{definition}[Overconfidence]
\label{def:overconf}
Overconfidence is the tendency to believe one's judgments are more accurate than they actually are. The literature distinguishes \emph{overprecision} (confidence in the accuracy of one's forecasts), \emph{overplacement} (belief that one is better than others), and \emph{overestimation} (unrealistically positive self-assessments of ability).
\end{definition}

The cleanest laboratory evidence is on overprecision. \citet{alpert1982} asked subjects to give 98\% confidence intervals for unknown quantities (the area of Manhattan, the number of piano tuners in Chicago). Across many studies, the realized coverage rate of these intervals is closer to 60\% than 98\%. Subjects are wildly miscalibrated: the ranges they assign are too narrow for the information they have.

\subsection{CFOs and the Stock Market}

The overprecision finding extends to professional investors. \citet{bendavid2013} surveyed chief financial officers of U.S.\ corporations and asked them to report probability distributions for the one-year-ahead return on the S\&P~500. The CFOs' distributions are very tight: their 80\% confidence intervals span about 28 percentage points of expected return, centered on a relatively narrow range of point estimates. The realized market return, over the same horizon, falls inside the CFOs' 80\% interval only about 36\% of the time. The professional respondents are not only miscalibrated, they are miscalibrated in the same direction as the student subjects of \citet{alpert1982}: their intervals are too narrow by a factor of two or more.

\begin{figure}[htbp]
\centering
\includegraphics[width=0.55\textwidth]{slide07_pic4.jpg}
\caption[CFO overprecision: subjective volatility too low]{\textbf{CFO overprecision.} Kernel densities of the one-year S\&P~500 return volatility implied by CFOs' survey responses, plotted against the historical distribution of realized one-year volatility over 1950--2000 and 2001--2011. The CFO-implied densities sit well inside the historical ones: executives' subjective return distributions are far too tight, which is why the realized return falls inside their 80\% confidence intervals only about 36\% of the time. \textit{Source:} \citet{bendavid2013}.}
\label{fig:cfo}
\end{figure}

\subsection{Overplacement}

A separate strand of evidence concerns \emph{overplacement}---the belief that one is above average. \citet{svenson1981} surveyed drivers in the U.S.\ and Sweden and found that 93\% of American respondents (and a large majority of Swedish respondents) rate themselves as above-average drivers. The figure is mechanically impossible: the average driver is, by definition, average. The same pattern appears in finance: in surveys of investor confidence, retail investors systematically report that their own portfolio will outperform the average portfolio of their peers.

\subsection{Excessive Trading}

Overconfidence has direct trading implications. \citet{odean1999} showed that the average household in a large discount-brokerage sample turns over roughly 80\% of its portfolio each year and that net returns fall sharply with turnover: the gross return on the portfolio is roughly uncorrelated with trading intensity, but trading costs---commissions and bid--ask spreads---rise with it. The households that trade most earn less than the households that trade least, and the average household earns less than a passive index fund. \citet{hong2007} document the matching aggregate fact that trading volume is far in excess of what any rational model with symmetric information can generate.

\begin{figure}[htbp]
\centering
\includegraphics[width=0.55\textwidth]{slide08_pic6.png}
\caption[Trading volume and disagreement]{\textbf{Excessive trading.} The figure shows that turnover in equity markets is far higher than would be implied by liquidity-motivated or rebalancing trades alone, and that the volume is concentrated in periods when investor beliefs diverge. The dispersion in beliefs is the natural reading of the volume pattern. \textit{Source:} EF8083 course materials, after \citet{hong2007}.}
\label{fig:excess_trading}
\end{figure}

\subsection{Gender Differences}

\citet{barber2001} show that the overconfidence bias is not gender-neutral. In their brokerage dataset, men trade roughly 1.5 times as often as women and earn correspondingly lower net returns; the gender gap in net returns is largest among singles and shrinks (but does not disappear) among married households, when household portfolios are pooled. \citet{lundeberg1994} document the matching psychological result: in laboratory studies of confidence judgments, men report higher confidence than women on tasks of comparable difficulty, with the gap widening on tasks at which both perform poorly. The two strands of evidence are mutually reinforcing: the gender gap in trading costs is consistent with the gender gap in confidence, and the mechanism operates through over-trading rather than through any systematic gender difference in stock-selection skill.

\begin{figure}[htbp]
\centering
\includegraphics[width=0.42\textwidth]{slide09_pic3.png}
\hspace{0.3em}

\includegraphics[width=0.42\textwidth]{slide09_pic4.png}
\caption[Gender, overconfidence, and trading]{\textbf{Boys will be boys.} Regression evidence from \citet{barber2001}: men, and single men in particular, trade at higher monthly turnover, earn lower net returns, and hold more volatile portfolios than women---the pattern predicted if men are the more overconfident sex \citep{lundeberg1994}. \textit{Source:} \citet{barber2001} and \citet{lundeberg1994}.}
\label{fig:gender}
\end{figure}

\subsection{Momentum and Overconfidence}

\citet{daniel1998} develop a formal model in which overconfidence interacts with \emph{self-attribution bias} to produce momentum. Investors are overconfident about private signals they themselves generate (e.g., their own research) but less so about public signals. Positive returns reinforce their overconfidence through self-attribution---a streak of good calls is attributed to skill, a streak of bad calls to bad luck---and they continue to trade on their private signals with growing positions. Eventually, the public information that contradicts their private view arrives, and prices partially correct. The dynamics produce short-horizon continuation (momentum) followed by partial reversal as information is revealed.

The model predicts that momentum should be stronger in environments where overconfidence is culturally or institutionally more pronounced. \citet{chui2010} test this prediction using Hofstede's individualism index as a proxy for cultural variation in overconfidence. Individualism, in Hofstede's framework, captures the extent to which a society rewards individual initiative and self-promotion relative to group cohesion---a pattern that maps naturally onto overplacement. The data confirm the prediction: over February 1984--June 2003, the winner-minus-loser momentum profit for country-average portfolios rises with individualism, from the bottom 30\% to the top 30\% of the individualism ranking, with a high-minus-low spread of 0.65\% per month ($t = 4.3$).

\begin{figure}[htbp]
\centering
\includegraphics[width=0.9\textwidth]{TableIV_PanelA_Momentum_and_Individualism.png}
\caption[Momentum profits by individualism]{\textbf{Momentum and individualism.} Winner and loser returns and winner-minus-loser momentum profits for country-average portfolios grouped by Hofstede's individualism index (bottom 30\%, medium, top 30\%), February 1984--June 2003. Momentum profits rise with individualism; the high-minus-low spread in the momentum profit is 0.65\% per month ($t = 4.30$). \textit{Source:} EF8083 course materials, based on \citet{chui2010}.}
\label{fig:ctw}
\end{figure}

% ============================================================
\section{Over-Extrapolation}
\label{sec:overextrapolation}

\begin{definition}[Over-extrapolation]
\label{def:overextrapolation}
The estimate of the future value of a quantity is a positive function of the recent past values of that quantity. The mechanism is the representativeness heuristic: a recent history of, say, increasing prices is judged representative of a generally increasing trend, so the next observation is expected to be larger still.
\end{definition}

Over-extrapolation is the most widely applied idea in behavioral finance. It is the source of return extrapolation (predicting future returns from past returns), of fundamentals extrapolation (predicting future earnings from past growth), and, in aggregate, of the survey-implied ``expected returns'' that move with valuation ratios in ways rational discount-rate models cannot easily reproduce. The applications span almost every asset pricing puzzle of the past twenty years.

\begin{figure}[htbp]
\centering
\includegraphics[width=0.65\textwidth]{13.png}
\caption[The over-extrapolation feedback loop]{\textbf{Over-extrapolation in pictures.} Good news lifts recent observations, which are extrapolated to lift beliefs about future values; the higher beliefs in turn push prices further up; the price rise confirms the recent observations. The loop runs until some shock breaks it. \textit{Source:} EF8083 course materials.}
\label{fig:loop}
\end{figure}

\subsection{Survey Evidence on Extrapolation}

\citet{greenwood2014} assemble survey expectations of one-year-ahead stock market returns from six sources, including the Gallup poll, the University of Michigan's survey of consumer sentiment, the Survey of Professional Forecasters, and several investor surveys. They compare the survey-implied expected return to realized subsequent returns over the same horizon and find a strongly negative correlation: the higher investors' expected returns, the lower realized returns are subsequently. The sign is the central fact of the paper. Survey expectations rise after market run-ups---they extrapolate recent returns---and high expectations therefore coincide with the conditions that precede low subsequent returns, the opposite of what rational time-varying-discount-rate models imply for measured expected returns. The authors then connect survey expectations to investor flows into equity mutual funds and show that the flows chase the same states of the world: inflows rise when expected returns are high, exactly the pattern that produces overvaluation when sentiment reverses.

\begin{figure}[htbp]
\centering
\includegraphics[width=0.55\textwidth]{slide14_pic4.png}
\caption[Survey expectations and realized returns]{\textbf{Survey expectations of stock returns.} Survey measures of one-year-ahead return expectations are negatively correlated with realized returns over the same horizon. The pattern is consistent with extrapolation: when recent returns are high, expected returns are also high, and subsequent returns are low---the expectations are inverted by the next period, not vindicated. \textit{Source:} \citet{greenwood2014}.}
\label{fig:gs_survey}
\end{figure}

\begin{figure}[htbp]
\centering
\includegraphics[width=0.55\textwidth]{slide15_pic4.jpg}
\caption[Survey expectations and investor flows]{\textbf{Investors act on extrapolative beliefs.} Net flows into equity mutual funds correlate with survey expectations of one-year-ahead returns. Investors buy after high expected returns (which follow high recent returns) and thereby deepen the overvaluation that the expectations forecast. \textit{Source:} \citet{greenwood2014}.}
\label{fig:gs_flows}
\end{figure}

\subsection{Measuring the Degree of Extrapolation: DOX}

\citet{cassella2018} propose a formal measure of the weight investors place on recent returns. Their measure, called \emph{DOX} for degree of extrapolation, is the complement of a parameter $\lambda$ in a forecaster's representation of next-period expected return as a weighted average of past returns,
\begin{equation}
\mathit{Exp}_t = a + b \sum_{i=0}^{\infty} w_i\, R_{t-(i+1)\Delta t,\; t-i\Delta t},
\qquad
w_i = \frac{\lambda^{i}}{\sum_{k=0}^{\infty} \lambda^{k}}, \quad 0 \le \lambda < 1,
\end{equation}
so that $\mathrm{DOX} = 1-\lambda$. A lower $\lambda$ places more weight on the most recent observations; earlier observations contribute less and less to the extrapolator's expectation. Estimating $\lambda$ recursively from survey data yields a time series for DOX rather than a constant: investors extrapolate more in some periods than in others (Figure~\ref{fig:dox}).

\begin{figure}[htbp]
\centering
\includegraphics[width=0.55\textwidth]{slide16_pic4.png}
\caption[DOX: degree of extrapolation]{\textbf{Degree of extrapolation.} The recursively estimated DOX time series for the principal component of the survey return expectations, 1992:1--2013:12. DOX varies over time: the weight investors place on past returns is a state variable, not a fixed parameter. \textit{Source:} \citet{cassella2018}.}
\label{fig:dox}
\end{figure}

DOX enters asset prices through the predictability of returns by price-scaled fundamentals (dividend yield, book-to-market, earnings-to-price). The predictability is contingent on the DOX. When DOX is high, investors more quickly ``forget'' the positive price changes that drove up valuations, and their beliefs mean-revert faster; as a result, high valuations are \emph{more} strongly associated with low future returns. When DOX is low, extrapolative beliefs persist longer and valuation ratios forecast less. The DOX measure is therefore not just a description of beliefs but a state variable that conditions the predictability of returns---and \citet{cassella2018} document that the conditional predictability is present in the data.

\begin{figure}[htbp]
\centering
\includegraphics[width=0.55\textwidth]{slide17_pic4.png}
\caption[Predictability of returns conditional on DOX]{\textbf{DOX and predictability.} Predictive regressions of subsequent returns on price-scaled variables, conditional on the estimated DOX measure. The slope is steeper (more predictability) when DOX is high: investors forget the price run-up faster, beliefs mean-revert faster, and high valuations more strongly forecast low future returns. \textit{Source:} \citet{cassella2018}.}
\label{fig:dox_cond}
\end{figure}

\subsection{Fundamentals Extrapolation}

A parallel stream of evidence concerns the extrapolation of fundamentals rather than prices. \citet{bordalo2023} show that long-term growth (LTG) expectations of corporate earnings overreact to good and bad news: a positive earnings surprise shifts the distribution of LTG forecasts upward by more than is justified by the signal alone. The authors track an aggregate LTG-expectation measure over time and show that its deviations from a fundamental benchmark are highly correlated with the price--dividend ratio and predict low subsequent market returns over a five-year horizon. The mechanism is the same as in return extrapolation: investors apply a diagnostic forecasting rule that overweights representative recent observations, and the implied LTG co-moves with prices and reverts.

\begin{figure}[htbp]
\centering
\includegraphics[width=\textwidth]{slide18_pic3.png}
\caption[LTG expectations, the price--dividend ratio, and returns]{\textbf{Fundamentals extrapolation.} The figure overlays a measure of long-term earnings-growth (LTG) expectations (red) with the price--dividend ratio (green) and five-year-ahead realized stock returns (black). LTG expectations track the price--dividend ratio and predict low subsequent returns, consistent with extrapolative beliefs. \textit{Source:} \citet{bordalo2023}.}
\label{fig:ltg}
\end{figure}

\subsection{My Research: Na\"ive Earnings Growth Extrapolation}
\label{sec:myresearch_ytd}

\citet{cuizhang2021} study fundamentals extrapolation in a setting where the extrapolated quantity is unusually visible: Chinese earnings announcements. China reports quarterly earnings on a year-to-date basis, which makes cumulative YTD growth highly salient to investors even though the YTD measure contains stale earnings information---earlier quarters' results roll mechanically into the cumulative number. Stocks with the strongest past YTD earnings growth run up in the five days before announcements and reverse afterward: investors extrapolate the salient YTD growth into the upcoming announcement and are disappointed when the ``growth'' turns out to be old news. The pattern concentrates in small, retail-heavy, low-analyst-coverage stocks, and retail buying demand rises in the run-up to announcements. The paper is the cross-sectional counterpart of \citet{bordalo2023} on the aggregate market: fundamentals extrapolation that distorts the prices of individual stocks around earnings news rather than the valuation of the market as a whole.

% ============================================================
\section{Other Biases: Conservatism, Anchoring, and Feelings}
\label{sec:otherbiases}

The biases catalogued so far all push investors to \emph{overreact} to new information---to be more confident in their forecasts, more certain about the trend, and more aggressive in their trading. The literature also identifies biases that push in the opposite direction, underreacting to news.

\subsection{Conservatism}

\begin{definition}[Conservatism]
\label{def:conservatism}
\emph{Conservatism} is the tendency to update beliefs too slowly in response to new information \citep{edwards1968}. In the canonical Bayesian account, agents place posterior probability somewhere between the prior and the signal-consistent likelihood; the conservatism bias is the finding that the updating weight on the signal is smaller than the Bayesian formula implies.
\end{definition}

The conservatism bias is the natural complement to over-extrapolation. Conservatism by itself would predict slow price adjustment to news and the post-earnings-announcement drift of Section~\ref{sec:pead}; over-extrapolation by itself would predict sharp initial reaction followed by reversal. \citet{barberis1998} build a model in which conservatism at short horizons and representativeness at long horizons combine to produce exactly the joint pattern of momentum and long-term reversal that the data show.

\subsection{Belief Perseverance and Confirmation Bias}

Two related biases affect how investors \emph{search} for and \emph{interpret} evidence. \emph{Belief perseverance} is the tendency, once an opinion is formed, to be slow to revise it on receipt of new evidence; subjects often continue to endorse a belief even after the evidence that produced it has been fully discredited. \emph{Confirmation bias} is the active counterpart: investors seek out evidence that supports their prior and interpret ambiguous evidence as supportive. The two effects compound: investors both fail to look for falsifying evidence and discount it when they encounter it. \citet{stanovich2013} show that higher cognitive ability can amplify confirmation bias, because more capable reasoners are better at constructing arguments for their existing position.

\subsection{Anchoring}

\begin{definition}[Anchoring]
\label{def:anchoring}
\emph{Anchoring} is the tendency, in forming numerical estimates of uncertain quantities, to adjust insufficiently from some initial value (the ``anchor''), even when the anchor is generated randomly and is clearly irrelevant \citep{kahneman1974}.
\end{definition}

In the canonical demonstration, \citet{kahneman1974} spin a wheel of fortune in front of subjects to produce a random number, ask subjects whether the percentage of African nations in the United Nations is greater or less than that number, and then ask them to estimate the actual percentage. The random anchor moves the estimates even though the subjects watch it being generated: those who compared their estimate to 10 subsequently estimated 25 percent, while those who compared it to 60 estimated 45 percent. The financial applications include IPO pricing (where the initial price range functions as an anchor for analysts), analyst price targets, and merger valuations (where the announced offer price anchors subsequent negotiation).

The application with the cleanest evidence is anchoring on the 52-week high. \citet{george2004} show that nearness to the 52-week high predicts future returns: a stock trading close to its high earns anomalously high subsequent returns, because traders are slow to react to good news when the price is already near a salient reference point. The same anchor contributes to the momentum effect---prices drift upward as the anchor itself rises.

\subsection{The Effect of Feelings}

A more recent strand of research investigates how feelings affect beliefs and preferences. \emph{Affective feelings}---mood and emotion---shift judgments in ways unrelated to the evidence: a happier mood produces more optimistic assessments, a sadder mood more pessimistic. \citet{edmans2006} document the case in finance: when the national soccer team loses a World Cup elimination match, the national stock market falls significantly the next trading day, with the effect concentrated in countries where soccer is more important. The economic mechanism is not information about fundamentals but mood: disappointing sports news worsens the national mood and depresses risk-taking.

Two environmental results reinforce the mood channel. \citet{hirshleifer2003} find higher stock market returns on sunnier days: sunshine improves mood, improved mood raises optimism about future cash flows, and the optimism is priced. \citet{kamstra2003} link the seasonal pattern in returns to Seasonal Affective Disorder: fewer hours of daylight induce a mild depression, which raises risk aversion and lowers returns through the winter months.

\citet{wang2020} extend the argument to a starker exogenous shock. Terrorist attacks raise aggregate risk aversion, as proxied by flows out of equity mutual funds and into government bond funds. The effect is larger in the days immediately following the attack and decays over time. Both findings point to a direct channel from affective feelings to beliefs about risk and, through beliefs, to portfolio choices. They also suggest a mechanism that can amplify market fluctuations: if the stock market rises, investors' mood improves, their judgments about the economy's prospects become more positive, and they push the market up still further.

\subsection{My Research: Air Pollution and Carbon Divestment}
\label{sec:myresearch_pollution}

Belief shifts triggered by exogenous stimuli reach professionals, and they can overshoot. \citet{huynh2025} study mutual fund managers exposed to local air pollution. Managers experiencing pollution underweight the stocks of high-emission firms, and the effect strengthens for environmentally conscious managers and for managers surprised by the pollution---that is, when the stimulus is salient, managers revise their beliefs about carbon risk. The divestment is costly: the carbon-intensive stocks managers sell subsequently outperform the stocks bought in their place. Like soccer losses and terrorist attacks, pollution is an exogenous stimulus that shifts beliefs; the difference is that here the professionals \emph{overreact}, selling carbon-intensive stocks by more than fundamentals justify.

\subsection{Cognitive Feelings}

\emph{Cognitive feelings} are about the ease with which information is processed rather than the mood that processing leaves behind. \citet{zajonc1968} and \citet{bornstein1989} document the basic result: people prefer stimuli that are familiar, easily processed, and fluently retrieved, even when the preference has no informational basis. \citet{uppal2003} formalize the asset pricing version: under ambiguity aversion or model misspecification, investors view stocks whose payoffs they can better estimate as less ambiguous, and overweight them. The applications are \emph{home bias} (domestic stocks are more familiar than foreign stocks) and \emph{local bias} (stocks headquartered nearby are more familiar than distant ones), both of which Chapter~\ref{ch:intro} documents.

% ============================================================
\section{Heterogeneous Beliefs: Sources}
\label{sec:hetsources}

The biases of Sections~\ref{sec:representativeness}--\ref{sec:otherbiases} all keep the \emph{representative investor} assumption intact: a single agent, possibly with systematically distorted beliefs, faces the asset pricing problem. The next two sections drop that assumption and work with heterogeneous beliefs---different agents holding different views about asset values, with the average belief unbiased. The departure is consequential. Asymmetric information alone cannot sustain rational disagreement: under common priors, a Bayesian agent who is offered a trade recognizes that the offer itself reflects information he does not have, and revises his beliefs until the apparent gains from trade vanish. The \citet{milgrom1981} no-trade theorem formalizes the result: with common priors and rational updating, agents who observe the same signals agree to agree, and no trade takes place. Heterogeneous beliefs in the rational sense must therefore come from somewhere beyond asymmetric information: from differences in prior beliefs, which \citet{morris1995} argues are reasonable for genuinely new phenomena, or from differences in how agents process the same information. The empirically important sources of such differences are the subject of this section.

The magnitude of the phenomenon is easy to underestimate. Figure~\ref{fig:spf} plots disagreement among respondents to the Survey of Professional Forecasters: the interquartile range across forecasters of one-quarter-ahead forecasts of annualized quarterly growth in real GDP, corporate profits, and CPI inflation. These are professionals, surveyed at the same moment, answering the same question with the same public data---and the disagreement is wide in every period. People regularly hold different beliefs about almost everything, from political views and sports events to economic outcomes and asset returns.

\begin{figure}[htbp]
\centering
\includegraphics[width=0.9\textwidth]{spf_dispersion.png}
\caption[Disagreement among professional forecasters]{\textbf{Even professionals disagree.} Interquartile range across respondents to the Survey of Professional Forecasters of one-quarter-ahead forecasts of annualized quarterly growth in real GDP, corporate profits, and CPI inflation (Philly Fed variables D2 and D1), 1968--2026. Forecasters answering the same question at the same time, with the same public data, disagree widely. \textit{Source:} Author's calculations based on the Federal Reserve Bank of Philadelphia Survey of Professional Forecasters.}
\label{fig:spf}
\end{figure}

\subsection{Experience Effects}

\begin{figure}[htbp]
\centering
\includegraphics[width=\textwidth]{slide30_pic3.jpg}
\caption[Personal experience and beliefs]{\textbf{Experience effects.} The figure illustrates how personal experience shapes beliefs: people place disproportionate weight on events they have lived through and form reference points from those events. The same statistical evidence can produce different posterior beliefs for agents with different life histories. \textit{Source:} EF8083 course materials.}
\label{fig:experience}
\end{figure}

One important source of heterogeneous beliefs is personal experience. When forming beliefs, a rational individual should take account of all available information; mounting evidence instead suggests that people overweight what they have lived through. \citet{malmendier2011} show that whether a household participates in the stock market, and how much it invests, depends on the stock market returns the household experienced over its lifetime: cohorts that lived through poor markets stay underweight equities for decades, and the effect of experienced returns fades only slowly as new experience accumulates.

\citet{malmendier2016} document a clean case in inflation expectations. Survey respondents overweight the inflation they experienced during their own lifetimes when forming expectations of future inflation: respondents born in the 1930s report higher inflation expectations in 2010 than respondents born in the 1970s, even though the cohort born later has experienced more inflation in absolute terms. Younger respondents update their beliefs more quickly in response to recent news; older respondents update more slowly and rely more heavily on their formative experiences. \citet{kuchler2019} extend the experience logic with individual-level data: people's beliefs about aggregate housing prices and unemployment depend strongly on the economic conditions they have experienced locally, even though the aggregate outcomes are public information. The same mechanism---overweighting personally experienced events---produces persistent cross-sectional dispersion in beliefs about any quantity for which the historical record is long.

\begin{figure}[htbp]
\centering
\includegraphics[width=\textwidth]{slide31_pic3.png}
\caption[Inflation expectations and lived experience]{\textbf{Inflation expectations.} Survey respondents' expectations of long-run inflation overweight the inflation experienced during their own lifetimes. The cohort born in the 1930s reports higher expected inflation in 2010 than the cohort born in the 1970s, opposite to what the realized record implies. \textit{Source:} \citet{malmendier2016}.}
\label{fig:malmendier}
\end{figure}

\subsection{Political Polarization}

\citet{goldman2025} document a more recent source of belief dispersion in financial news. The authors compare the same corporate financial news covered by the conservative \emph{Wall Street Journal} and the liberal \emph{New York Times}. The two newspapers use systematically more positive language for firms aligned with their political leanings, and the gap is wider for politically sensitive topics (climate, labor, executive compensation) than for neutral ones. Investors who consume one source rather than the other form different beliefs about the same companies, and the dispersion persists because the consumers do not switch sources. The finding is part of a broader pattern: media fragmentation along political lines is a quantitatively significant source of disagreement about financial facts.

\begin{figure}[htbp]
\centering
\includegraphics[width=\textwidth]{slide32_pic4.png}
\caption[Political polarization in financial news]{\textbf{Political polarization in financial news.} Tone-of-language scores for corporate financial coverage by the \emph{Wall Street Journal} (conservative) and \emph{New York Times} (liberal), separated by firm-level political alignment. Both newspapers use systematically more positive language for firms aligned with their political leanings, with the gap larger for politically sensitive topics. \textit{Source:} \citet{goldman2025}.}
\label{fig:polar}
\end{figure}

\subsection{Overconfidence in Private Signals}

The third source is the one emphasized by \citet{daniel1998}. Two agents with the same prior can observe the same public signal and apply different likelihoods because they differ in their confidence about their own private information. The agent with higher confidence in his private signal will place less weight on the public signal and more on his own view; the agent with lower confidence will do the reverse. Both are acting rationally, given their subjective confidence, but they hold heterogeneous beliefs about the asset, and the dispersion in beliefs is larger when the agents are overconfident in their private signals. The mechanism is the bridge from the overconfidence of Section~\ref{sec:overconfidence} to the heterogeneous-beliefs models of the next section.

% ============================================================
\section{Heterogeneous Beliefs with Short-Sale Constraints}
\label{sec:hetwithssc}

Heterogeneous beliefs alone do not have interesting pricing effects. If two agents disagree about an asset's value, the optimist wants to buy and the pessimist wants to sell. Without frictions, the two positions can be unwound and the price lies at a weighted average of the two beliefs. The overpricing and underpricing cancel out, and the equilibrium expected return is the rational expected return under the average belief. What dispersion does generate is volume: investors trade because they disagree about asset values, and higher dispersion of beliefs means higher trading volume. \citet{grinblatt2009}, using Finnish data that link brokerage records to aptitude tests, show that more overconfident individuals---those whose self-reported confidence exceeds what their test performance justifies---trade more. Heterogeneous beliefs become interesting for \emph{prices} when pessimists cannot fully express their view---in the canonical case, because short-selling is constrained.

\subsection{Static Argument: Miller (1977)}

\citet{miller1977} observed that if pessimists face binding short-sale constraints, only the optimists' valuation is reflected in the price. The equilibrium price lies above the average belief and above the asset's fundamental value, with the magnitude of the overpricing determined by the dispersion of beliefs. The implication is a sharp empirical prediction: assets that investors disagree about more should be more overpriced, and should earn lower subsequent returns.

The argument is static, in the sense that it does not require beliefs to change over time. It is enough that beliefs differ across agents in the cross-section at a given date, and that the pessimists cannot fully short. The mechanism is also asymmetric across the two legs of the trade: correcting overpricing requires shorting, which is constrained, while correcting underpricing requires only cash. The asymmetry is a deep source of pricing distortion in the heterogeneous-beliefs model.

\begin{figure}[htbp]
\centering
\includegraphics[width=0.55\textwidth]{34.png}
\caption[Heterogeneous beliefs and overpricing with short-sale constraints]{\textbf{Miller (1977) in pictures.} The figure illustrates the static argument: with short-sale constraints, only the optimist's valuation is reflected in the price, and the equilibrium price exceeds the average belief. The magnitude of the overpricing scales with the dispersion of beliefs. \textit{Source:} EF8083 course materials.}
\label{fig:miller}
\end{figure}

\subsection{Dynamic Argument: Harrison and Kreps (1978)}

The dynamic version of the argument adds a \emph{resale option}. \citet{harrison1978} show that an agent with optimistic beliefs about an asset is willing to pay more than his buy-and-hold valuation, because there is a positive probability that future information will make his view the prevailing one, at which point he can resell to a future optimist at a still higher price. The resale option is valuable precisely because beliefs are heterogeneous: in a homogeneous-belief world there is no one to resell to. With short-sale constraints preventing pessimists from driving the price down, the resale option adds a \emph{bubble component} to the price that can persist even when all agents know the asset is overvalued relative to fundamental value.

The following example, drawn from the \citet{harrison1978} framework, illustrates the magnitude of the resale-option effect (Figure~\ref{fig:resale}).

\begin{figure}[htbp]
\centering
\includegraphics[width=0.7\textwidth]{resale_option.png}
\caption[The resale-option example]{\textbf{The resale-option example.} Valuations and prices in the two-agent, two-state example. Both agents' buy-and-hold valuations at $t=0$ are 50, yet both are willing to pay 57.5, because each expects to resell to the more optimistic agent at $t=1$: at 90 in state $u$, above $B$'s valuation of 75, or at 25 in state $d$, above $A$'s valuation of 10. \textit{Source:} EF8083 course materials.}
\label{fig:resale}
\end{figure}

\begin{quote}
\textbf{Worked example.} Two risk-neutral agents $A$ and $B$ trade a risky asset at $t = 0, 1, 2$. The interest rate is zero and short-selling is prohibited. At $t = 1$ a public signal arrives---good news (state $u$) or bad news (state $d$), each with probability $1/2$---and the two agents interpret it asymmetrically. After good news, agent $A$ values the asset at 90 and agent $B$ at 75; after bad news, agent $A$ values it at 10 and agent $B$ at 25. Each agent's buy-and-hold valuation at $t = 0$ is therefore $(90+10)/2 = 50$ for $A$ and $(75+25)/2 = 50$ for $B$: without the option to resell, the two agents agree the asset is worth 50. With short-selling prohibited, the $t = 1$ price is set by the more optimistic agent: 90 in state $u$ and 25 in state $d$. At $t = 0$, both agents are therefore willing to pay
\[
\tfrac{1}{2} \cdot 90 + \tfrac{1}{2} \cdot 25 \;=\; 57.5,
\]
\emph{above} either agent's buy-and-hold valuation. The excess 7.5 is the resale-option premium. Each agent expects to resell the asset to the other at a price above his own valuation: $B$ expects to resell to $A$ at 90 in state $u$, which exceeds his valuation of 75, and $A$ expects to resell to $B$ at 25 in state $d$, which exceeds his valuation of 10. The resale option enters both agents' initial willingness to pay: a speculative bubble component.
\end{quote}

\subsection{Bubble Characteristics}

The bubble component of \citet{harrison1978} has properties that match observed bubble episodes. \citet{scheinkman2003} formalize the result. First, the resale-option premium is determined by fluctuations in agents' beliefs: the more likely beliefs are to flip between optimists and pessimists, the more the agents trade and the more valuable the resale option---which is why bubbles tend to coincide with trading frenzies. Second, the premium depends on asset float: with more shares available for trading, the bubble is smaller. The mechanism runs through downward-sloping demand curves. An owner can sell his share only if a buyer's belief passes his own by a sufficient margin; a larger float makes it less likely that enough sufficiently optimistic buyers exist, and the resale option is worth less. The predictions are exactly the dimensions on which \citet{xiong2011} take the model to Chinese warrant data (Section~\ref{sec:hetprices}).

% ============================================================
\section{Heterogeneous Beliefs and Asset Prices}
\label{sec:hetprices}

\subsection{High Valuations and Heavy Trading}

Models of heterogeneous beliefs with short-sale constraints are popular partly because they explain not just overpricing but also \emph{the coincidence of high valuations and heavy trading}, which is one of the empirical regularities of bubble episodes (Chapter~\ref{ch:intro}). With short-sale constraints, the optimists trade against each other as their beliefs diverge; pessimists who would otherwise trade to express their views are sidelined. The trading volume and the price level both rise, and they do so together. \citet{hong2007} document this co-movement in U.S.\ stock market data.

\begin{figure}[htbp]
\centering
\includegraphics[width=0.40\textwidth]{slide39_pic4.jpg}
\caption[High valuations and trading volume]{\textbf{High valuations and heavy trading.} The figure plots turnover against cumulative returns in the cross-section of U.S.\ stocks. Stocks with the highest turnover also exhibit the highest cumulative returns, consistent with heterogeneous-beliefs models in which optimists trade against each other while pessimists are constrained. \textit{Source:} \citet{hong2007}, reproduced from the original article.}
\label{fig:valtrading}
\end{figure}

\subsection{The Chinese Warrant Bubble}

The most rigorous empirical test of the heterogeneous-beliefs-with-short-sale-constraints theory is \citet{xiong2011} on Chinese warrants. Between 2005 and 2008, more than a dozen put warrants traded on the Chinese exchanges at intrinsic values that were essentially zero (the underlying stocks were deep out of the money relative to exercise prices). Yet the warrants traded at high prices and with enormous turnover. The fundamental value of a deep-out-of-the-money put warrant is very small; the observed prices are an order of magnitude larger.

\begin{figure}[htbp]
\centering
\includegraphics[width=\textwidth]{slide40_pic3.png}
\caption[Chinese warrant prices and turnover]{\textbf{Chinese warrant bubble.} The figure shows price and turnover for a representative deep-out-of-the-money put warrant traded in China, 2005--2008. Intrinsic value is essentially zero throughout the sample; the warrant trades at high prices and with very high turnover. \textit{Source:} \citet{xiong2011}, reproduced from the original article.}
\label{fig:warrant}
\end{figure}

\citet{xiong2011} document three regularities consistent with the resale-option theory of bubbles. The size of the bubble is increasing in turnover (consistent with high dispersion of beliefs driving both prices and volume); increasing in the volatility of the underlying asset (consistent with the resale-option value rising in the volatility of beliefs); and decreasing in the asset float (consistent with tighter short-sale constraints when fewer shares are available). The warrant data are an unusually clean test because the fundamental value is essentially known and constant, leaving little room for the rational discount-rate counter-story.

\begin{figure}[htbp]
\centering
\includegraphics[width=0.55\textwidth]{slide41_pic4.png}
\caption[Bubble size and its determinants]{\textbf{Warrant bubble evidence.} The bubble component of warrant prices increases with turnover and volatility and decreases with asset float, consistent with the resale-option theory of speculative bubbles. \textit{Source:} \citet{xiong2011}.}
\label{fig:warrant_determinants}
\end{figure}

\subsection{Cross-Sectional Stock Returns}

The heterogeneous-beliefs model with short-sale constraints also predicts a cross-sectional return pattern: assets that investors disagree about more should be more overpriced, and should earn lower subsequent returns. \citet{diether2003} test the prediction using analyst forecast dispersion as the disagreement measure. They form portfolios sorted on the cross-sectional standard deviation of analyst earnings forecasts and find that high-dispersion stocks earn significantly lower returns than low-dispersion stocks over the subsequent twelve months, even after controlling for size, book-to-market, and momentum. The return spread is consistent with the overpricing mechanism: disagreement is high precisely for stocks where some analysts are optimistic and others are pessimistic, and the short-sale constraint prevents the pessimists from fully expressing their view. Figure~\ref{fig:dms} shows the pattern in the Open Source Asset Pricing data: over 1976--2024, the low-dispersion quintile out-earns the high-dispersion quintile by a statistically significant spread.

\begin{figure}[htbp]
\centering
\includegraphics[width=0.7\textwidth]{ForecastDispersion_quintile_bar.png}
\caption[Analyst disagreement and subsequent returns]{\textbf{Cross-sectional disagreement.} Average monthly value-weighted returns (\%) of quintile portfolios sorted on analyst earnings-forecast dispersion, 1976--2024; quintile 1 holds the highest-dispersion stocks. Low-dispersion stocks outperform high-dispersion stocks (spread $t = 3.2$), consistent with the overpricing prediction of the heterogeneous-beliefs model with short-sale constraints. \textit{Source:} Author's calculations based on Open Source Asset Pricing quintile portfolio data \citep{chen2022}, after \citet{diether2003}.}
\label{fig:dms}
\end{figure}

\subsection{Aggregate Stock Market}

The same mechanism predicts a relation at the aggregate level: when aggregate disagreement is high, the market as a whole should be more overpriced and earn lower subsequent returns. \citet{yu2011} measures aggregate disagreement as the value-weighted average of stock-level forecast dispersion and shows that high-dispersion periods are followed by low returns on the aggregate market index.

\begin{figure}[htbp]
\centering
\includegraphics[width=0.45\textwidth]{slide43_pic5.png}
\hspace{0.3em}
\includegraphics[width=0.45\textwidth]{slide43_pic4.png}
\caption[Aggregate disagreement and market returns]{\textbf{Aggregate disagreement.} Left: the time series of aggregate analyst forecast dispersion, measured as the value-weighted average of stock-level dispersion. Right: a nonparametric estimate of expected market returns against aggregate disagreement, with pointwise 95\% confidence bands; expected returns decline as disagreement rises. High-disagreement periods are followed by low market returns, consistent with the heterogeneous-beliefs overpricing mechanism operating at the aggregate level. \textit{Source:} \citet{yu2011}, reproduced from the original article.}
\label{fig:yu}
\end{figure}

\subsection{Risk--Return Tradeoff}

Another important application of the same logic is the risk--return tradeoff. Standard theory predicts that higher risk is rewarded with higher expected return: \citet{sharpe1964} and the CAPM that follows from it imply an upward-sloping security market line. Empirically, the relation between risk and return is often negative instead. Heterogeneous-beliefs models explain the anomaly: when investors disagree about a stock, the optimists set the price and the pessimists are sidelined by the short-sale constraint, so higher disagreement is associated with a \emph{higher} price and a \emph{lower} expected return. The negative risk--return relation is one of the empirical regularities of Chapter~\ref{ch:intro}. Figure~\ref{fig:alphabeta} shows the pattern in beta-sorted portfolios: average CAPM alphas decline with beta across the ten BetaFP decile portfolios, with cross-sectional fit $\alpha = 1.178 - 0.178\,\beta$ and $R^2 = 0.95$. High-beta stocks do not earn commensurately high returns; measured against the CAPM, they are overpriced.

\begin{figure}[htbp]
\centering
\includegraphics[width=0.6\textwidth]{beta.png}
\caption[CAPM alphas decline with beta]{\textbf{Alpha versus beta.} Average monthly CAPM alphas against CAPM betas for the ten beta-sorted (BetaFP) decile portfolios, Open Source Asset Pricing data, 1928--2024 ($N = 1{,}161$ months; market proxy: equal-weighted mean of the deciles, no risk-free rate available). The fitted line $\alpha = 1.178 - 0.178\,\beta$ ($R^2 = 0.95$) slopes down: high-beta portfolios earn lower alphas, the high-beta--low-alpha pattern that standard theory cannot produce. \textit{Source:} Author's calculations based on Open Source Asset Pricing decile portfolio data \citep{chen2022}.}
\label{fig:alphabeta}
\end{figure}

\citet{hong2016} turn the observation into a model of the \emph{high-beta--low-alpha puzzle}. Write each asset's payoff as its beta times the aggregate cash flow plus idiosyncratic noise,
\[
d_i = \beta_i(Z) + e_i .
\]
Suppose agents disagree about the aggregate cash flow $Z$: optimists expect $Z+\delta$, pessimists $Z-\delta$. The disagreement passes through beta. High-beta assets are more sensitive to disagreement about aggregate cash flows than low-beta assets, so optimists expect $\beta_i(Z+\delta)$ and pessimists $\beta_i(Z-\delta)$: high-beta assets experience more divergence of opinion. With short-sale constraints, the amplified divergence overprices high-beta assets relative to low-beta assets. In a slogan, \emph{beta amplifies aggregate disagreement}.

The model delivers sharp predictions about the shape of the security market line. When aggregate disagreement is low, the CAPM holds and the SML slopes upward. When aggregate disagreement is high, the SML takes an inverted-U shape---expected returns initially increase with beta and then decrease as overpricing of high-beta assets dominates. Stock-level disagreement should increase with beta, especially when aggregate disagreement is high. \citet{hong2016} measure aggregate disagreement as the lagged-beta-weighted average of stock-level analyst forecast dispersion (Figure~\ref{fig:hs_fig1}) and test the predictions in beta-decile portfolios split by high and low disagreement months (Figure~\ref{fig:specbeta}).

\begin{figure}[htbp]
\centering
\includegraphics[width=0.5\textwidth]{slide45_pic2.png}
\caption[Aggregate disagreement, Hong and Sraer measure]{\textbf{Aggregate disagreement.} The time series of aggregate disagreement from \citet{hong2016}: each month, the lagged-beta-weighted average across stocks of the standard deviation of analyst forecasts of the stock's long-run EPS growth, December 1981--February 2014, CRSP stocks excluding penny stocks and microcaps. \textit{Source:} \citet{hong2016}, reproduced from the original article.}
\label{fig:hs_fig1}
\end{figure}

\begin{figure}[htbp]
\centering
\includegraphics[width=0.48\textwidth]{slide46_pic3.png}
\hspace{0.3em}
\includegraphics[width=0.48\textwidth]{slide46_pic4.jpg}
\caption[Speculative betas in the data]{\textbf{Speculative betas in the data.} Left: average three-month excess returns of equal-weighted beta-decile portfolios during low and high aggregate-disagreement months. Right: the equal-weighted average dispersion of analyst earnings forecasts by beta decile during low and high aggregate-disagreement months. The panels are the empirical counterpart of the model's predictions that stock-level disagreement rises with beta and that the shape of the security market line depends on aggregate disagreement. \textit{Source:} \citet{hong2016}.}
\label{fig:specbeta}
\end{figure}

\subsection{Macro Factors and Disagreement}

\citet{li2016} extends the logic from individual stocks to the pricing of macroeconomic factors. If investors disagree about the prospects of consumption growth or inflation, the resulting overpricing should be more pronounced for factor portfolios whose payoffs load more heavily on the most-disputed macro states. The data are consistent with the prediction: factor returns are lower when disagreement about the underlying macro state is higher, and the effect is stronger for factor portfolios more exposed to that state.

\begin{figure}[htbp]
\centering
\includegraphics[width=0.55\textwidth]{slide47_pic4.png}
\caption[Macro factor disagreement]{\textbf{Macro factors and disagreement.} The figure shows that the returns on factor portfolios more exposed to highly disputed macro states are lower in periods of high disagreement about those states. \textit{Source:} \citet{li2016}, reproduced from the original article.}
\label{fig:li_macro}
\end{figure}

% ============================================================
\section{Causal Evidence: The EDGAR Identification}
\label{sec:causal}

The cross-sectional and time-series correlations of Section~\ref{sec:hetprices} are consistent with the heterogeneous-beliefs model, but the same correlations admit other readings. A stock with high analyst dispersion is also a stock with high uncertainty, high trading volume, and high media attention; the overpricing interpretation must compete with risk-based and information-based alternatives. Distinguishing the channels requires exogenous variation in disagreement.

\citet{chang2022} provide such variation by exploiting the staggered rollout of the SEC's EDGAR electronic disclosure system. Before EDGAR, investors could access firms' mandatory filings only at high cost, either by subscribing to commercial data providers or by physically visiting the SEC's reference rooms. Beginning in April 1993, the SEC required U.S.\ firms to file their mandatory disclosures electronically through the EDGAR system, and it randomly assigned firms to one of ten implementation waves, staggering inclusion over the three-year period between 1993 and 1996. The rollout therefore moves information-acquisition costs exogenously across firms and time. The authors argue that lower information costs reduce the dispersion of beliefs: as more investors access the same public information, disagreement about what the firm is worth shrinks. EDGAR inclusion helps resolve disagreement around information events, leading to price corrections and lower crash risk.

\begin{figure}[htbp]
\centering
\includegraphics[width=0.40\textwidth]{slide49_pic4.png}
\hspace{0.2em}
\includegraphics[width=0.40\textwidth]{slide49_pic5.png}
\caption[EDGAR rollout and disagreement]{\textbf{Causal evidence on disagreement.} Evidence from the EDGAR quasi-experiment of \citet{chang2022}: as staggered EDGAR inclusion lowers information-acquisition costs, disagreement resolves around information events, prices correct, and crash risk falls (right: crash-risk regressions using NCSKEW). \textit{Source:} \citet{chang2022}, reproduced from the original article.}
\label{fig:chang}
\end{figure}

The central finding is that the rollout of EDGAR, by lowering information costs and disagreement, also lowered stock prices---the sign predicted by the overpricing interpretation. The magnitude is consistent with a quantitatively meaningful share of the variation in prices being driven by disagreement, and the design rules out the alternative interpretations that rely on time-invariant firm characteristics. The result is one of the cleanest causal pieces of evidence in the heterogeneous-beliefs literature, and it sharpens the case for treating disagreement as a priced source of risk in its own right rather than as a proxy for some other dimension of firm fundamentals.

% ============================================================
\section{Do People Ever Learn?}
\label{sec:stakes}

The biases of Sections~\ref{sec:representativeness}--\ref{sec:otherbiases} are persistent; they are part of the cognitive and emotional machinery that evolution built into human decision-making. Many of the most persistent behavioral biases are mental shortcuts with survival value: loss aversion increases survival probability, and overconfidence increases effort and trust within one's community. The kinds of financial decisions in which the biases matter most are recent on evolutionary timescales. Do people simply learn their way out of them? Only to a limited extent.

Do highly trained professionals exhibit the biases? Yes, and sometimes professionals do worse. CEOs and founders are more overconfident than the average person---selection into those roles rewards confidence, and success at the top of the distribution reinforces it. Expertise does not immunize judgment; it supplies the wrong mental models with more resources.

The most direct evidence compares investor behavior under different stakes. \citet{sui2025} use a unique setting in which investors trade stocks in real accounts with their own money and simultaneously in a simulated setting where the personal stakes are lower. Within-investor estimation shows that the same investors exhibit stronger biases---disposition effect, lottery preference, extrapolation, underdiversification, local bias---and perform worse in the higher-stakes real accounts than in the simulated accounts. Investors do better when they care less. The explanation is uncomfortable for the debiasing view: individuals carry incorrect mental models, and greater attention and effort in the high-stakes setting lead to the more careful \emph{application} of those incorrect models. Trying harder makes the biases worse, not better.

\begin{figure}[htbp]
\centering
\includegraphics[width=0.7\textwidth]{slide51_pic4.png}
\caption[Stakes and investor biases]{\textbf{Investors do better when they care less.} For each bias---disposition effect, lottery preference, extrapolation, underdiversification, local bias---the table reports the mean across investors' real accounts, the mean in their simulated accounts, and the within-investor difference with $t$-statistics and Wilcoxon $p$-values ($N$: number of investors). Biases are systematically stronger in the real, higher-stakes accounts. \textit{Source:} \citet{sui2025}, reproduced from the original article.}
\label{fig:sui}
\end{figure}

The finding has two implications for the behavioral asset pricing program. First, the biases are not diluted by stakes; they are magnified. The emotional salience of real money leads investors to apply their (incorrect) mental models more carefully, so raising the stakes deepens rather than corrects the distortion. Second, it raises a methodological point that recurs in the chapter: behavioral biases are not fixed properties of the investor. They are properties of the investor \emph{in context}, and the context matters for how strongly they affect trading and prices.

% ============================================================
\section{What Remains Contested}
\label{sec:beliefscontested}

Two strands of evidence keep the behavioral reading honest.

First, the cognitive biases catalogued in Sections~\ref{sec:representativeness}--\ref{sec:otherbiases} are individually well-documented in laboratory data, but their \emph{quantitative} contribution to asset prices is harder to pin down. The disposition effect of Chapter~\ref{ch:preferences} and the trading-cost evidence of Section~\ref{sec:overconfidence} are reasonably direct: investors' biases translate into overtrading and into measurable return patterns. The over-extrapolation evidence is stronger at the survey level than at the trading level: the Greenwood--Shleifer correlations are tight, but the contribution of extrapolative beliefs to the cross-section of stock returns is harder to estimate cleanly. The bias-by-bias literature must be aggregated across biases that may or may not interact, and the empirical case is uneven.

Second, the heterogeneous-beliefs models with short-sale constraints are competitive with rational explanations in several applications. The cross-sectional disagreement result of \citet{diether2003} admits a risk-based interpretation: high-dispersion stocks are also stocks with high fundamental uncertainty, and the low returns may be compensation for that uncertainty rather than the result of overpricing. The warrant-bubble evidence of \citet{xiong2011} is sharper, but it remains a single market; the macro-factor extension of \citet{li2016} is suggestive rather than conclusive. The causal EDGAR evidence of \citet{chang2022} is the strongest single piece of evidence, but it speaks to one channel (information costs and disagreement) rather than to the broader class of bias-driven models.

The methodological lesson is the same as in the earlier chapters: a behavioral model earns credibility when it predicts something not used to fit it, and when that prediction survives out-of-sample tests. The strongest out-of-sample predictions of the biased-beliefs program are the disposition effect and the trading-cost evidence of overconfidence, the Greenwood--Shleifer correlations of survey expectations and returns, the warrant-bubble evidence on resale-option values, and the EDGAR causal evidence on disagreement. The rest of the literature is competitive with rational alternatives.

% ============================================================
\section{The Road Ahead}
\label{sec:beliefsahead}

This chapter has developed the second of the two departures from the rational benchmark---biased beliefs---together with the cross-sectional dispersion in beliefs that the heterogeneous-beliefs model exploits. Combined with the prospect-theory departures of Chapter~\ref{ch:preferences}, the book now has the two psychological ingredients needed for most of behavioral finance: a non-standard theory of preferences under risk, and a non-standard theory of belief formation. The combination is what most applications in the chapters that follow deploy.

The remaining chapters assemble these ingredients into specific empirical and theoretical applications. Chapter~\ref{ch:bounded} develops \emph{bounded rationality}: limited attention, category-based thinking, and psychology-free models of noise trading and sentiment. A further chapter develops \emph{memory-based} approaches---the way investors form beliefs from the cases they remember rather than from a Bayesian aggregation of evidence. These models retain rational beliefs in a thin sense but relax the assumption that all available information is processed.

Together with Chapter~\ref{ch:intro} on the empirical facts, Chapter~\ref{ch:limits} on limits to arbitrage, Chapter~\ref{ch:preferences} on preferences, and the present chapter on biased beliefs, the book develops the building blocks of behavioral asset pricing. The remaining chapters apply them to specific settings---the role of attention, the formation of investor sentiment, and the pricing of memory-driven beliefs---in each case tracing the path from a documented cognitive bias to a measurable asset-pricing implication.
